\chapter{Projects, packages, tests, and automation}
\label{ch:packages-tests}

\begin{learningobjectives}
After this chapter, you should be able to:
\begin{itemize}
  \item distinguish scripts, modules, import packages, and distributions;
  \item explain dependency declarations, lockfiles, and version identifiers;
  \item select unit, integration, contract, system, and regression tests from a
  concrete risk; and
  \item distinguish continuous integration, delivery, and deployment.
\end{itemize}
\end{learningobjectives}

\begin{chapterconnection}
A well-designed function or class is useful in one notebook. A package makes it
available across notebooks, scripts, collaborators, and versions. That larger
audience creates compatibility and change risks, so packaging, testing, and
continuous integration belong in the same engineering story.
\end{chapterconnection}

\section{From exploration to reusable software}

Notebooks are valuable for combining narrative, code, and evidence. They become
fragile when hidden execution order, copied definitions, mutable global state,
and machine-specific paths accumulate. Mature behavior should move into a
module where notebooks and tests import the same implementation.

A \term{script} is a file intended to be executed as a program. A
\term{module} is an importable Python file. An \term{import package} organizes
modules under a shared namespace. A \term{distribution} is an installable
project released through a packaging mechanism; its distribution name need not
equal its import name.

\begin{center}
\begin{tikzpicture}[node distance=7mm]
  \node[wideflowbox] (notebook) {Notebook\\question and evidence};
  \node[wideflowbox,right=of notebook] (module) {Package module\\reusable behavior};
  \node[wideflowbox,right=of module] (tests) {Tests\\executable contracts};
  \draw[flowarrow] (notebook) -- node[above,font=\small]{refactor} (module);
  \draw[flowarrow] (tests) -- node[above,font=\small]{verify} (module);
  \draw[flowarrow] (module) to[bend left=25] node[below,font=\small]{import} (notebook);
\end{tikzpicture}
\end{center}

The \code{src} layout places import packages below \code{src/}. It helps detect
accidental imports from the repository root: tests exercise the installed
project instead of a convenient shadow copy.

\section{A public package API}

A package should deliberately expose names that callers may rely on and keep
implementation helpers private. Re-exporting selected objects from
\code{rice\_dsm/\_\_init\_\_.py} creates a facade that can remain stable when
internal modules move.

Every public name is a compatibility promise. Its behavior needs documentation,
types, error contracts, examples, and tests. ``Public'' means supported; Python
may still technically allow importing a private underscore-prefixed name.

\section{Dependencies and versions}

\code{pyproject.toml} declares project metadata and dependency constraints.
\code{uv.lock} records an exact resolved graph for reproducible setup. The local
environment installs that graph. These artifacts answer different questions:

\begin{tabularx}{\textwidth}{@{}p{0.25\textwidth}X@{}}
\toprule
Artifact & Question \\
\midrule
Dependency constraint & Which releases are intended to be compatible? \\
Lockfile & Which complete graph did this project resolve? \\
Installed environment & Which files are available to this interpreter now? \\
Package version & Which release of our public behavior is this? \\
Code revision & Which exact source state produced this artifact? \\
\bottomrule
\end{tabularx}

A version number is useful only when tied to an artifact and policy. A Git
commit identifies source precisely; a semantic version communicates intended
compatibility; a data version identifies evidence; a model version identifies a
trained artifact. Do not substitute one for another.

\section{Tests begin with risks}

A test should protect a meaningful behavior or failure mode. Start with a risk,
choose the smallest boundary that can reveal it, and use an oracle independent
enough to detect the defect.

\begin{longtable}{@{}p{0.18\textwidth}p{0.29\textwidth}p{0.43\textwidth}@{}}
\toprule
Test kind & Boundary & Typical purpose \\
\midrule
\endhead
Unit & One function or class & Domain cases, validation, and local invariants \\
Property & Many generated values & General algebraic or structural rules \\
Metamorphic & Related executions & Relations known even when exact output is hard \\
Integration & Multiple real components & Database, filesystem, or package interaction \\
Contract & Producer and consumer agreement & Schema and compatibility without full deployment \\
System & Whole running application & Behavior across assembled components \\
Regression & Previously observed failure & Prevent return of a specific defect \\
Data quality & Dataset and pipeline & Schema, ranges, uniqueness, drift, and leakage signals \\
\bottomrule
\end{longtable}

Passing tests show consistency with encoded expectations. They do not prove
that the expectations are scientifically sufficient. Test doubles are useful
when they isolate one contract; excessive mocking can produce a fictional
system that never tests the real dependency.

\section{Automation and CI/CD}

\term{continuous integration} (CI) automatically checks integrated changes in
a clean environment. \term{continuous delivery} keeps a validated release
deployable through an explicit decision. \term{continuous deployment}
automatically releases every qualifying change to production.

They are not synonyms. A course repository needs CI to prevent broken notebooks
or package behavior from entering the protected branch. It does not need
automatic production deployment to teach that boundary.

\begin{center}
\begin{tikzpicture}[node distance=5mm]
  \node[flowbox] (change) {Change};
  \node[flowbox,right=of change] (static) {Lint and\\build};
  \node[flowbox,right=of static] (test) {Tests and\\notebooks};
  \node[flowbox,right=of test] (review) {Review and\\required gate};
  \draw[flowarrow] (change) -- (static);
  \draw[flowarrow] (static) -- (test);
  \draw[flowarrow] (test) -- (review);
\end{tikzpicture}
\end{center}

The course workflow synchronizes the committed lockfile, builds the
distribution, verifies the package source and kernel, runs tests, and executes
every notebook on Linux, macOS, and Windows. Cross-platform CI catches path,
shell, filesystem, and environment assumptions that one laptop cannot reveal.

\begin{professionalbox}
CI is evidence, not responsibility. A green check cannot detect an untested
requirement, choose a valid statistic, approve a data-use policy, or monitor a
deployed client. Branch protection makes the check consequential by preventing
unverified changes from entering the protected branch.
\end{professionalbox}

\begin{checkpointbox}
For each risk, choose a test boundary: a probability function mishandles $p=0$;
a SQLite transaction fails to roll back; a browser client expects a removed
field; a notebook imports a package absent from the lockfile; and a model
pipeline leaks future labels.
\end{checkpointbox}

\section{Worked example: derive tests from a risk}

Suppose \code{shortest\_path(source, target)} returns a sequence of graph node
identifiers. A weak test might assert one result from one graph. A stronger test
plan begins by asking how the behavior could be wrong.

\begin{workedexample}{protect shortest-path behavior}
\textbf{Risk 1: the function returns a path containing a nonexistent edge.}
For every adjacent pair in the result, assert that the graph contains that edge.
This is a structural property and does not require knowing the exact path in
advance.

\textbf{Risk 2: the path is valid but not shortest.} On a small fixture, compare
its length with a simple independent breadth-first search oracle or with
manually established distances.

\textbf{Risk 3: source equals target.} Specify whether the result is a
single-node path. Test it explicitly; do not expect a random graph fixture to
exercise the case.

\textbf{Risk 4: no path exists.} Decide whether the API returns \code{None}, an
empty tuple, or raises a domain exception. Encode one stable behavior.

\textbf{Risk 5: unknown identifiers.} Verify that source and target errors are
distinguishable and do not silently resemble a disconnected graph.

\textbf{Risk 6: input mutation.} Snapshot graph state before the call and assert
that a read-only query does not alter nodes or edges.
\end{workedexample}

Different tests can cover different boundaries. Unit tests protect graph
semantics in memory. An integration test can load the graph from CSV and JSON
first. A command-line test can verify human-facing output and exit status. A
notebook execution test can prove that the instructional path still imports the
packaged implementation.

\section{Fixtures, parameters, and doubles}

A \term{fixture} establishes controlled test state and owns cleanup. Small
fixtures should make the important relationship visible: a line graph, cycle,
disconnected pair, or graph with two equally short routes.

Parametrization applies the same behavioral claim to a table of cases. It is
most useful when each case has the same meaning and failure message. A long
table of unrelated cases can become harder to diagnose than separate tests.

A test double replaces a dependency to isolate one contract. For example, a
fake clock can make time deterministic and a stubbed cloud response can test
parsing without credentials. But a fake database cannot prove that actual SQL,
transactions, constraints, or driver behavior work. The test boundary must
match the claim.

\begin{misconceptionbox}
\textbf{Misconception:} a unit test should mock every collaborator.

\textbf{Correction:} a unit is a coherent behavior, not necessarily one
function isolated from all real objects. Mock only a boundary whose behavior is
outside the claim or costly and unsafe to invoke. Prefer simple real value
objects when they make the test more faithful.
\end{misconceptionbox}

\section{Reading a CI failure as evidence}

A CI job runs commands in a new machine context. A failure that appears only on
Windows may reveal separator, case, shell, path-length, or file-locking
assumptions. A failure that appears only from a clean checkout may reveal an
uncommitted local file or hidden notebook state. A lockfile failure may reveal
that dependency declarations changed without resolution.

Read the first causal error, record the operating system and command, reproduce
the smallest relevant boundary, then add a regression check. Rerunning until a
flaky check passes discards evidence and permits uncertainty into the protected
branch.

\conceptcheck{A test passes locally but fails in CI because \code{data.csv}
cannot be found. List three distinct hypotheses and one observation that would
separate them.}

\section{Exercises}

\subsection*{Foundational}
\begin{enumerate}
  \item Distinguish a module, import package, and distribution using the
  \code{rice-dsm} project.
  \item Explain why a lower bound such as \code{numpy>=2} and a lockfile answer
  different version questions.
  \item Contrast continuous integration, continuous delivery, and continuous
  deployment without using the word ``automatic'' as the entire definition.
\end{enumerate}

\subsection*{Applied}
\begin{enumerate}
  \item Choose one function in \code{rice\_dsm}. Write a risk inventory and map
  each risk to a test kind and independent oracle.
  \item Convert three repetitive edge-case tests into a clear parametrized test.
  Explain what would make parametrization inappropriate.
  \item Read the course CI workflow. For each command, state the failure class it
  is intended to reveal and one class it cannot reveal.
\end{enumerate}

\subsection*{Design}
Propose a release gate for a small scientific package. Include source checks,
package build, unit and integration tests, notebook or example execution,
artifact identity, review, and a response to flaky tests.

\begin{chapterreview}
\textbf{Key ideas.} Notebooks explore and explain; modules preserve reusable
behavior. A public package API is a compatibility promise. Declarations,
lockfiles, installed environments, package versions, data versions, and code
revisions identify different artifacts. Tests arise from risks. CI makes
selected checks repeatable in clean environments.

\textbf{Vocabulary.} script; module; import package; distribution; public API;
dependency; version constraint; lockfile; fixture; parametrization; test double;
oracle; unit; property; integration; contract; system; regression; CI/CD.

\textbf{Explain aloud.} Why is ``all tests passed'' weaker than ``the scientific
result is valid''?
\end{chapterreview}

\begin{notebooklink}
Lecture 3 notebooks 00 through 02 build the project mental model, refactor a
knowledge graph into the package, and develop a detailed testing and automation
taxonomy.
\end{notebooklink}

\section{Takeaway}

Packaging creates a maintained interface across files and users. Tests encode
selected promises. Automation makes those checks repeatable and consequential.
Together they allow scientific software to change without abandoning evidence.
